% This LaTeX document needs to be compiled with XeLaTeX.
\documentclass[10pt]{article}
\usepackage[utf8]{inputenc}
\usepackage{amsmath}
\usepackage{amsfonts}
\usepackage{amssymb}
\usepackage[version=4]{mhchem}
\usepackage{stmaryrd}
\usepackage{graphicx}
\usepackage[export]{adjustbox}
\graphicspath{ {./images/} }
\usepackage[fallback]{xeCJK}
\usepackage{polyglossia}
\usepackage{fontspec}
\usepackage{newunicodechar}
\IfFontExistsTF{Noto Serif CJK JP}
{\setCJKmainfont{Noto Serif CJK JP}}
{\IfFontExistsTF{STSong}
  {\setCJKmainfont{STSong}}
  {\IfFontExistsTF{Droid Sans Fallback}
    {\setCJKmainfont{Droid Sans Fallback}}
    {\setCJKmainfont{SimSun}}
}}

\setmainlanguage{english}
\IfFontExistsTF{CMU Serif}
{\setmainfont{CMU Serif}}
{\IfFontExistsTF{DejaVu Sans}
  {\setmainfont{DejaVu Sans}}
  {\setmainfont{Georgia}}
}

\newunicodechar{✓}{\ifmmode\checkmark\else{$\checkmark$}\fi}

\begin{document}
\section*{Asymptotic decay rate of an infinite product of sinc functions}
Asked 8 years, 4 months ago Modified 7 years, 7 months ago Viewed 1k times

\begin{itemize}
  \item[-] (related to my previous question)\\
20 Consider the function
  \item[✓] 
\begin{equation*}
f(x)=\prod_{n=0}^{\infty} \operatorname{sinc}\left(\frac{\pi x}{2^{n}}\right), \quad x \geq 0, \tag{1}
\end{equation*}

  \item[ロ] where $\operatorname{sinc}(z)$ denotes the sinc function. The function $f(x)$ has zeros at positive integers, and oscillates with a quickly decaying amplitude. Its signs on the intervals between consecutive zeros follow the same pattern as the Thue-Morse sequence.\\
\includegraphics[max width=\textwidth, alt={}, center]{c02ddc48-cf27-411a-9a03-1741e9e77c8a-1_1001_1637_1250_286}
\end{itemize}

Can we find a real-valued function $g(x)$ (elementary, if possible), such that it is analytic, positive and monotone decreasing (and having monotone derivatives of any order, if possible) and satisfies


\begin{equation*}
\lim _{t \rightarrow \infty} \frac{1}{t-t_{0}} \int_{t_{0}}^{t} \frac{|f(x)|}{g(x)} d x=1 \tag{2}
\end{equation*}


for large enough $t_{0}$ ? Or, at least,


\begin{equation*}
\exists A>0, \exists B>A, \forall t>t_{0}, A<\frac{1}{t-t_{0}} \int_{t_{0}}^{t} \frac{|f(x)|}{g(x)} d x<B \tag{3}
\end{equation*}


for large enough $t_{0}$ ?\\
The "for large enough" provision means that $g(x)$ is permitted to be not monotone, or have zeros or discontinuities, or be not defined at all for small $t$, and we only care about its "eventual" behavior.

Empirically, it looks like $g(x)$ should decay faster than any negative power of $x$, but slower than exponentially. I'm thinking something close to $\exp \left(-\log ^{2} x\right)$, but perhaps not exactly that.

Update: Maybe this expansion can be useful:


\begin{equation*}
\prod_{n=0}^{m} \operatorname{sinc}\left(\frac{\pi x}{2^{n}}\right)=\frac{2^{\binom{m}{2}}}{(\pi x)^{m+1}} \sum_{n=0}^{2^{m}-1} t_{n} \sin \left(\frac{\pi m}{2}+\frac{2 n+1}{2^{m}} \pi x\right), \tag{4}
\end{equation*}


where $t_{0}=1, t_{n}=(-1)^{n} t_{\lfloor n / 2\rfloor}$ (the signed Thue-Morse sequence).

Formulas for $\boldsymbol{\operatorname { s i n c }}(\boldsymbol{\pi} \mathbf{z})$

Using symmetry formulas for the gamma function


\begin{equation*}
\Gamma(z) \Gamma(1-z)=\frac{\pi}{\sin (\pi z)} \tag{2}
\end{equation*}


easy to obtain


\begin{gather*}
\frac{\sin (\pi z)}{\pi z}=\frac{1}{z \Gamma(z) \Gamma(1-z)}, \\
\operatorname{sinc}(\pi \mathrm{z})=\frac{1}{\Gamma(1+z) \Gamma(1-z)},  \tag{3}\\
\ln \operatorname{sinc}(\pi z)=-(\ln \Gamma(1+z)+\ln \Gamma(1-z)) . \tag{4}
\end{gather*}


Using the Maclaurin series in the form of


\begin{equation*}
\ln \Gamma(1+z)=-\ln (1+z)+(1-\gamma) z+\sum_{j=2}^{\infty}(-1)^{j} \frac{\zeta(j)-1}{j} z^{j}, \quad|z|<2 \tag{5}
\end{equation*}


where $\gamma=0.57721 \ldots$, one can get

$$
\begin{gathered}
\ln \Gamma(1-z)=-\ln (1-z)-(1-\gamma) z+\sum_{j=2}^{\infty} \frac{\zeta(j)-1}{j} z^{n} \\
\ln \operatorname{sinc}(\pi z)=\ln \left(1-z^{2}\right)-\sum_{k=1}^{\infty} \frac{\zeta(2 k)-1}{k} z^{2 k}
\end{gathered}
$$

Zeta function of the even argument is


\begin{equation*}
\zeta(2 k)=\frac{(2 \pi)^{2 k}}{2(2 k)!}\left|B_{2 k}\right|, \tag{6}
\end{equation*}


where $B_{2 k}$ are the Bernoulli numbers. Then


\begin{equation*}
\ln \operatorname{sinc}(\pi z)=\ln \left(1-z^{2}\right)-\sum_{k=1}^{\infty} c_{k} z^{2 k}, \tag{7}
\end{equation*}


where


\begin{equation*}
c_{k}=\frac{1}{2 k}\left(\frac{(2 \pi)^{2 k}}{(2 k)!}\left|B_{2 k}\right|-2\right) . \tag{8}
\end{equation*}


Formulas for $\mathbf{R}(\mathbf{x})$

Let us find the sum

$$
\begin{gathered}
\ln R_{m}(x)=\sum_{n=m}^{\infty} \ln \operatorname{sinc} \frac{\pi x}{2^{n}} \\
\ln R_{m}(x)=\sum_{n=m}^{\infty}\left(\ln \left(1-\left(\frac{x}{2^{n}}\right)^{2}\right)-\sum_{k=1}^{\infty} c_{k}\left(\frac{x}{2^{n}}\right)^{2 k}\right)= \\
\sum_{n=m}^{\infty} \ln \left(1-\left(\frac{x}{2^{n}}\right)^{2}\right)-\sum_{k=1}^{\infty}\left(c_{k} x^{2 k} \sum_{n=m}^{\infty} 2^{-2 n k}\right)= \\
\sum_{n=m}^{\infty} \ln \left(1-\left(\frac{x}{2^{n}}\right)^{2}\right)-\sum_{k=1}^{\infty} \frac{c_{k}}{1-2^{-2 k}}\left(\frac{x}{2^{m}}\right)^{2 k} .
\end{gathered}
$$

The first term can be presented over the q-Pochhammer Symbol in the form of


\begin{equation*}
\sum_{n=0}^{m} \ln \left(1-\left(\frac{x}{2^{n}}\right)^{2}\right)=\ln \left(x^{2}, \frac{1}{4}\right)_{m} . \tag{9}
\end{equation*}


This allows to obtain the formulas


\begin{equation*}
R(x)=\prod_{n=0}^{m-1} \operatorname{sinc} \frac{\pi x}{2^{n}} \cdot \frac{\left(x^{2}, \frac{1}{4}\right)_{\infty}}{\left(x^{2}, \frac{1}{4}\right)_{m}} \cdot \exp \left(\sum_{k=1}^{\infty} \frac{c_{k}}{1-2^{-m k}}\left(\frac{x}{2^{m}}\right)^{2 k}\right), \quad|x|<2^{m+1} . \tag{(10}
\end{equation*}


✓ C

Radius of convergence in the formula (10) shows that it gives $R(x)=0$ outside this radius, such as the Stirling asymptotics. Easy to see that parameter $m$ allows to manage this radius. At the same time, it allows to provide the reqired length of sum.

If the bounds of the required interval are predefined, then paraneter $m$ can be considered as constant. This allows to use the common function for this interval.

For example, for the predefined interval $x=(10,100)$ the value $m=6$ can be chosen. Then formula (10) presents the function $|R(x)|$ as the production of the three factors:\\
partial sinc production,\\
\includegraphics[max width=\textwidth, alt={}, center]{c02ddc48-cf27-411a-9a03-1741e9e77c8a-5_488_795_120_322}

\begin{verbatim}
plotprod_n = 0^5 | ((pi*x*2^(-n)) sin(pi*x*2^(-n))/
    ((pi*x*2^(-n))^2+10^-10)) |.
wherexin(10,100)
Computed by Wolfram|Alpha
\end{verbatim}

q-Pochhammer ratio\\
\includegraphics[max width=\textwidth, alt={}, center]{c02ddc48-cf27-411a-9a03-1741e9e77c8a-5_539_791_988_322}

\begin{verbatim}
plot(prod_(n=5)^20(1.0-x^22^(-2*n-2)))*
    sign(64-x), wherexin(10,100)
Computed by Wolfram|Alpha
\end{verbatim}

and the reminder series exponent\\
\includegraphics[max width=\textwidth, alt={}, center]{c02ddc48-cf27-411a-9a03-1741e9e77c8a-5_539_791_1865_322}

\begin{verbatim}
plotprod_n =
    l^3e^((|bernou|li(2n)|/(2n)!(2pi)^(2n)-2)/
        (1-64^(-n))
        (x/64)^(2n)/(2n)), wherex=10, 100
Computed by Wolfram|Alpha
\end{verbatim}

In the general case, the most convenient way looks using of formula (10) for a piecewise approximation on the required intervals .

Anyway, this allows to take in account the multiplicity of preceeding zeros.

\begin{verbatim}
Share Cite Edit Follow Flag edited Jan 24, 2019 at 19:52
\end{verbatim}

\begin{center}

\includegraphics[max width=\textwidth, alt={}]{c02ddc48-cf27-411a-9a03-1741e9e77c8a-6_142_428_395_1446}
\end{center}

Anyway, the function $f(x)$ is known to be continuous and analytic, its Taylor-Maclaurin series and Fourier transform are studied pretty well, see, for example, the references in my previous question.

\begin{itemize}
  \item[-] Vladimir Reshetnikov Apr 30, 2018 at 3:56
\end{itemize}

\begin{itemize}
  \item @VladimirReshetnikov New version. Waiting for comments. - Yuri Negometyanov May 1, 2018 at 5:06
\end{itemize}

In formula (10), there is a free variable $m$ on the right, but not on the left -- is it a typo? Also, the summation over $k$ in (10) starts from $k=0$, but the summand contains $c_{k}$. Looking at the definition (8), we can see that $c_{0}$ is undefined because it involves a division by zero. Also, there is another division by zero because of the denominator $1-2^{-m k}$.- Vladimir Reshetnikov May 1, 2018 at 23:10

\begin{itemize}
  \item @VladimirReshetnikov Of course, this was a typo. The earlier formulas contains $k=1$. Thank you for\\
accuracy. - Yuri Negometyanov May 1, 2018 at 23:20
  \item @VladimirReshetnikov By the way. Is known the inequaliity\\
$\frac{2(2 n)!}{(2 \pi)^{2 n}}\left(\frac{1}{1-2^{1-2 n}}\right)>(-1)^{n+1} B_{2 n}>\frac{2(2 n)!}{(2 \pi)^{2 n}}, \quad n=1,2, \ldots$. - Yuri Negometyanov May 1, 2018 at 23:36
\end{itemize}

What about the free variable $m$ in (10)? - Vladimir Reshetnikov May 2, 2018 at 0:26\\
I wonder where the scope of the $\prod$ in (10) ends. Is it over only the first factor, or all three?

\begin{itemize}
  \item Vladimir Reshetnikov May 2, 2018 at 0:30\\
@ - VladimirReshetnikov Really, (10) contains the ratio of $R(x)=R_{\infty}(x)$ and $R_{m}(x)$, and that allows
  \item to acheve the nesessary (arbitrary) accuracy. Taking in account the convergency raduis and the required possibly closed form, it does not looking so free. - Yuri Negometyanov May 2, 2018 at 6:32
  \item @VladimirReshetnikov It is possible to bound the required series length in (10). At the same time, it
  \item gives o if $x \geq 2^{m+1}$. And I don't see any way to get the asymptotic decay rate as the universal function in the closed form. - Yuri Negometyanov May 2, 2018 at 6:52
\end{itemize}

Thank you for the effort (and the bounty is yours), but I do not think this answer brings us significantly closer to resolving my question. - Vladimir Reshetnikov May 3, 2018 at 18:19

\begin{itemize}
  \item @VladimirReshetnikov The problem is solved in closed form for the 'local infinity" settings (see the\\
updated text). I think this is the maximum of possible. - Yuri Negometyanov May 4, 2018 at 1:40
\end{itemize}

\end{document}